% Generated draft: requires mathematical and source review.
\documentclass[11pt]{article}
\usepackage[margin=25mm]{geometry}
\usepackage{fontspec}
\setmainfont{Latin Modern Roman}
\usepackage{amsmath,amssymb,mathtools,unicode-math}
\setmathfont{Latin Modern Math}
\usepackage{enumitem,needspace,xurl,hyperref}
\setlength{\parindent}{0pt}
\setlength{\parskip}{5pt}
\setlength{\emergencystretch}{3em}
\newcommand{\N}{\mathbb{N}}
\newcommand{\Z}{\mathbb{Z}}
\newcommand{\Q}{\mathbb{Q}}
\newcommand{\R}{\mathbb{R}}
\newcommand{\C}{\mathbb{C}}
\newcommand{\F}{\mathbb{F}}
\newcommand{\A}{\mathbb{A}}
\newcommand{\PP}{\mathbb{P}}
\newcommand{\E}{\mathbb{E}}
\newcommand{\eps}{\varepsilon}
\newcommand{\dd}{\,\mathrm{d}}
\begin{document}
\section*{UnsolvedMath ID: 8800053}
\textbf{Source problem code:} catalogue entry
\subsection{Definitions and mathematical statement}
Let \(E/\F_q\) be an elliptic curve and let \(r\nmid q\) be a large prime dividing its order. Write
\[
\#E(\F_q)=h r,
\qquad k=\operatorname{ord}_r(q)
=\min\{j\ge1:r\mid q^j-1\}.
\]
Here \(h\) is the cofactor and \(k\) the embedding degree for the chosen subgroup order \(r\). The ratio \(\rho=\log q/\log r\) is another usual efficiency parameter, but no target \(\rho\) is stated in the source.

Lange's programme asks for additional MNT/MNT-like pairing-friendly constructions: more families, usable larger \(k\), and smaller cofactors under the intended parameter tradeoffs. A family must produce actual finite fields and elliptic curves with the indicated order/divisibility and embedding-degree conditions. Polynomial parameter identities alone do not prove infinitely many simultaneous prime parameter values. Usability requires specified subgroup/field sizes and computational criteria; it is not a mathematical bound supplied by the word usable.

The classical strict MNT terminology is associated with prime-order constructions at \(k=3,4,6\). In a prime-order setting \(h=1\), already the smallest positive cofactor. Thus a request for larger degrees and smaller cofactors must either concern a broader MNT-like class or compare different non-prime-order constructions; the exact intended class/baseline is missing. No new degree threshold or universal cofactor optimum is invented.

The supplied research summary reports substantial families, mentioning MNT, BN at \(k=12\), BLS, Freeman at \(k=10\), and KSS. It describes the 2006 programme as largely achieved for many degrees, not as a complete classification or an optimum for every degree/size. Those reported successes do not supply exact guarantees for every named family.
\subsection{Short English statement}
Construct further pairing-friendly elliptic-curve families with explicitly controlled embedding degrees and cofactors. Clarify whether MNT is used strictly or for a broader class, since prime-order MNT curves already have cofactor one.
\subsection{Sources}
\begin{itemize}
\item Local UnsolvedMath registry, permanent ID \(8800053\), exact source code \texttt{catalogue entry}, ``More MNT and pairing-friendly elliptic curves.'' Proposer Tanja Lange.
\item IPAM \emph{Workshop on Number Theory and Cryptography --- Open Problems} (2006), attachment 6239, recorded PDF page 41, printed slide 39; \url{https://www.ipam.ucla.edu/programs/workshops/workshop-i-number-theory-and-cryptography-open-problems/?tab=schedule}. Recorded access 2026-07-29; pdf-and-legacy-ppt visual-audit extraction.
\item Supplied PARTIAL-PROGRESS research summary names Miyaji--Nakabayashi--Takano (2001), Barreto--Naehrig (2005), Barreto--Lynn--Scott, Freeman, and KSS, without complete publication references or precise parameter/cost hypotheses.
\end{itemize}
\end{document}
